\documentclass{beamer}

\usepackage[ngerman]{babel}
\usepackage{graphicx}


% Example here:
% https://github.com/josephwright/beamer/blob/main/doc/solutions/conference-talks/conference-ornate-20min.de.tex

\mode<presentation>
{
	\usetheme{Singapore}
	%\usetheme{AnnArbor}
	%\usecolortheme{apricot}
}

\title{Kolloquium Pflichtenheft}
\subtitle{Speaker: Johann Schicho}

%\titlegraphic{\includegraphics[width=5cm]{../../docs/Pflichtenheft/graphics/LasEs-logo}}

\date{\small 2. November 2021}
\logo{\includegraphics[width=1.5cm]{../../docs/Pflichtenheft/graphics/LasEs-logo}}

\author{\textbf{Team 2} \\ \small {Johannes Garstenauer, Stefanie Gürster, Thomas Kirz,\\ Johann Schicho, Sebastian Vogt} \normalsize}


\begin{document}

	\begin{frame}
		\titlepage
	\end{frame}

	\begin{frame}{Gliederung}
		\tableofcontents
	\end{frame}

\section{Anforderungen und Funktionalität}

\begin{frame}{Anforderungen}
	\begin{itemize}
		\item \emph{Submission und Review Management Systems}
		\pause % Durchklicken wie bei Kraftpunkt.

		\item Aufteilung in Rollen
		\begin{itemize}
			\item Nutzer
			\item Gutachter
			\item Editor
			\item Administrator
		\end{itemize}
		\pause

		\item Java mit Jakarta Server Faces

	\end{itemize}
\end{frame}

\begin{frame}{Funktionalität}
	content...
\end{frame}

\end{document}